\begin{problem}[Part III synthesis]\label{prob:vi17-p21}
Explain why the pair $(\Spec A,\OO_{\Spec A})$ has local rings at every point and why this is exactly the data needed before defining an affine scheme.  Do not use the affine-scheme universal property.
\end{problem}

\ifdefined\FullProblemDossiers
\paragraph{Why this problem matters.}
The problem marks the conceptual boundary between sheaf theory and scheme theory without importing results reserved for later corpus Problems.

\paragraph{Useful ideas.}
Combine the stalk formula with the local-ring criterion; distinguish construction from the naming/geometry of affine schemes.

\paragraph{Guided hints.}
\begin{enumerate}
\item At $\mfp$ the stalk is $A_{\mfp}$.
\item Its maximal ideal is $\mfp A_{\mfp}$.
\item A space with a sheaf of rings having local stalks is a locally ringed space.
\end{enumerate}

\begin{solution}
The structure sheaf makes $\Spec A$ a ringed space.  At every prime $\mfp$, its stalk is $A_{\mfp}$, which is a local ring with unique maximal ideal $\mfp A_{\mfp}$.  Therefore $(\Spec A,\OO_{\Spec A})$ is a locally ringed space.  VI/18 calls this particular locally ringed space the affine scheme associated to $A$.  The definition of a general scheme as a locally affine locally ringed space is deferred to VI/21.  No morphism classification or basic-open subscheme isomorphism is needed for the present conclusion.
\end{solution}

\paragraph{What happened mathematically.}
All the local algebra needed for affine geometry has now been assembled; the next chapter changes viewpoint from construction to geometric objects.

\paragraph{Extensions.}
In VI/18 prove that every $D(f)$ with the restricted sheaf is itself an affine scheme.

\paragraph{Related problems.}
The exercises and advanced examples provide computational variants of this result.

\paragraph{Provenance.}
\textbf{New transition problem.  AG3 Problems 11--14 remain reserved for later chapters and are not duplicated here.}
\else
\paragraph{Provenance.} \textbf{New transition problem.  AG3 Problems 11--14 remain reserved for later chapters and are not duplicated here.}
\fi
